\documentclass{article}
\usepackage[UTF8]{ctex}
\usepackage{amsmath,amssymb,amsfonts}
\usepackage{enumitem}
\usepackage{geometry}
\geometry{a4paper, margin=1in}
% === 错题本排版优化（不参与题目标识） ===
\setlength{\parindent}{0pt}
\setlength{\parskip}{0.35em plus 0.10em minus 0.08em}
\setlength{\emergencystretch}{3em}
\setlist{itemsep=0.25em, topsep=0.40em, parsep=0pt}
\allowdisplaybreaks[2]
% === 错题本排版优化结束 ===
\begin{document}
\section*{复旦大学2026年数学分析真题}
\section*{35．复旦大学2026年分析试卷}

\subsection{题目1}
1．已知函数项级数 \(\sum_{n=2}^{\infty} \frac{\cos (n \pi x)}{n^{a x}}\) 关于 \(x \in(2,+\infty)\) 一致收敛，求实数 \(a\) 的取值范围．

\subsection{题目2}
2．计算积分 \(\int_{0}^{+\infty}\left(\frac{x^{2}}{1+x^{4}}\right)^{2} \mathrm{~d} x\) ．

\subsection{题目3}
3．求曲面 \(z=x y\) 被柱面 \(x^{2}+y^{2}=1\) 所截部分的面积．

\subsection{题目4}
4．已知椭圆 \(C=\left\{(x, y) \mid x^{2}+4 y^{2}=16\right\}\) ，方向取逆时针，计算曲线积分

\[
\int_{C}\left(\frac{x y^{2}}{x^{2}+2 y^{2}+1}-y\right) \mathrm{d} x+\left(x+\frac{x^{2} y}{x^{2}+2 y^{2}+1}\right) \mathrm{d} y .
\]

\subsection{题目5}
5．求 \(\sin \ln x\) 在 \(x=1\) 处带 Peano 余项的四阶 Taylor 展开式。

\subsection{题目6}
6．已知 \(f(x) \in C[0,1]\) 在 \((0,1)\) 上二阶可导，\(f(0)=f(1)=0, f\left(\frac{1}{2}\right)=\frac{1}{4}\) ，证明：存在 \(\xi \in(0,1)\) ，使得

\[
f^{\prime \prime}(\xi)=-2
\]

\subsection{题目7}
7．已知 \(\zeta(x)=\sum_{n=1}^{\infty} \frac{1}{n^{x}}\) ，证明：\(\zeta(x)\) 在 \((1,+\infty)\) 上任意阶导数存在，并判断 \(\int_{1}^{+\infty}(\zeta(x)-1) \mathrm{d} x\) 的敛散性。

\subsection{题目8}
8．计算极限

\[
\lim _{\varepsilon \rightarrow 0^{+}} \iint_{x^{2}+x y+y^{2} \geq \varepsilon^{2}} \frac{1}{x^{2}+y^{2}} \sin \frac{1}{x^{2}+y^{2}} \mathrm{~d} x \mathrm{~d} y
\]

并说明计算过程的合理性。

\subsection{题目9}
9．证明幂级数 \(\sum_{n=0}^{\infty} \frac{x^{3 n}}{(3 n)!}\) 的和函数是微分方程 \(\frac{\mathrm{d}^{2} y}{\mathrm{~d} x^{2}}+\frac{\mathrm{d} y}{\mathrm{~d} x}+y=e^{x}\) 的一个解，并求出和函数的具体表达式。

\subsection{题目10}
10．证明：\((0,0)\) 是自洽系统 \(\left\{\begin{array}{l}\frac{\mathrm{d} x}{\mathrm{~d} t}=y-f(x) \\ \frac{\mathrm{d} y}{\mathrm{~d} t}=-g(x)\end{array}\right.\) 的奇点，并判断其李雅普诺夫意义下的稳定性，其中\(f, g\) 是实值连续可微函数，满足 \(x \neq 0\) 时，\(x f(x)>0, x g(x)>0\) ．

\subsection{题目11}
11．设 \(f(z)\) 在区域 \(\Omega\) 上全纯，满足 \(\sin (f(z))=f(z)\) ，证明：\(f(z)\) 是常值函数。

\subsection{题目12}
12．设 \(f(z)\) 在复平面上全纯，且对任意 \(\operatorname{Re} z \neq 0\) 的 \(z \in \mathbb{C}\) 都有 \(|f(z)|<\frac{1}{\mid \operatorname{Re} z}\) 。证明：\(f(z)=0\) 。

\subsection{题目13}
13．设 \(f, g\) 是实直线 \(\mathbb{R}\) 上的勒贝格可积函数，证明：二元函数 \(F(x, y)=f(x) g(y+x)\) 是 \(\mathbb{R}^{2}\) 上的勒贝格可积函数。

\subsection{题目14}
14．证明黎曼勒贝格引理：对于 \([a, b]\) 上的勒贝格可积函数 \(f\) 以及 \([0,1]\) 上的勒贝格可积的周期函数 \(g\) ，有

\[
\lim _{n \rightarrow \infty} \int_{a}^{b} f(x) g(n x) \mathrm{d} x=\int_{0}^{1} g(x) \int_{a}^{b} f(x) \mathrm{d} x
\]
\end{document}
